\begin{recetaConImagen}{Tacos Estilo Pastor}{2 personas}{40 min activos · 4 h 40 min total}{receta:tacos_al_pastor}
    {imagenes/tacos_al_pastor.png}{Piña, cebolla y un toque de limón para rematar unos tacos bien sabrosos.}

    \begin{minipage}[t]{0.35\textwidth}
        \section*{Ingredientes}
        \begin{itemize}[leftmargin=*]
            \item 600 g de carne de cerdo en filetes finos
            \item 3 cucharadas de pimentón de la Vera
            \item 1 cucharadita de comino
            \item 1 cucharadita de cebolla en polvo
            \item 1/2 cucharadita de ajo y perejil
            \item 1/2 cucharadita de pimienta
            \item zumo de 2 naranjas
            \item 1/4 de taza de vinagre
            \item 300 g de piña
            \item 1 cebolla
            \item cilantro fresco
            \item 12 tortillas de maíz
            \item aceite y sal
        \end{itemize}
    \end{minipage}
    \hfill
    \begin{minipage}[t]{0.6\textwidth}
        \section*{Preparación}
        \begin{enumerate}
            \item \textbf{Preparar el adobo:} Mezclar el pimentón de la Vera, el comino, la cebolla en polvo, el ajo y perejil, la pimienta, el zumo de naranja y el vinagre hasta que quede bien integrado.
            \item \textbf{Marinar la carne:} Cubrir bien la carne de cerdo con el adobo y dejarla marinar durante 4 horas para que coja todo el sabor.
            \item \textbf{Empezar la cocción:} Hacer la carne a fuego medio-alto con un poco de aceite hasta que esté un poco hecha.
            \item \textbf{Añadir la piña y la cebolla:} Cortar la piña y la cebolla finas, incorporarlas a la sartén con la carne y seguir cocinando hasta que quede todo bien hecho y sabroso.
            \item \textbf{Picar:} Cortar la carne en tiras o trozos pequeños y reservar caliente.
            \item \textbf{Preparar el acompañamiento:} Picar el cilantro fresco.
            \item \textbf{Montar:} Calentar las tortillas y rellenarlas con la carne junto con la piña y la cebolla cocinadas. Añadir por encima el cilantro.
            \item \textbf{Servir:} Terminar cada tortita con un chorreón de limón y servir recién hechos.
        \end{enumerate}
    \end{minipage}

    \vspace{1em}
    \begin{tcolorbox}[colback=orange!5, colframe=orange!50, title=Nota, size=small]
        Esta es mi versión: el equilibrio entre el adobo, la piña salteada con la cebolla y el toque final de limón es lo que hace que estos tacos queden tan sabrosos.
    \end{tcolorbox}
\end{recetaConImagen}
